\documentclass[12pt,letterpaper]{article}

% Configuración de página y formato
\usepackage[left=2.4cm,right=2.4cm,top=2.4cm,bottom=2.4cm]{geometry}
\usepackage[utf8]{inputenc}
\usepackage[T1]{fontenc}
\usepackage[spanish]{babel}
\usepackage{mathptmx} % lo mismo que o mas similar a Times New Roman
\usepackage{setspace}
\usepackage{hyperref}
\usepackage{enumitem}

\setlength{\parindent}{0pt} % Sin sangría
\singlespacing % Renglón cerrado
\pagestyle{empty} 

\hypersetup{
    colorlinks=true,
    urlcolor=black,
    linkcolor=black,
    citecolor=black
}

% para subtítulos negrita, izquierda, sin punto final
\newcommand{\subtitulo}[1]{\noindent\textbf{#1}\par\vspace{2pt}}

\begin{document}


\begin{center}
\textbf{BUENAS PRACTICAS PARA REALIZAR ENCUESTAS}\\[8pt]
\textit{R. A. Ramírez Gómez}\\
7690-22-12700 Universidad Mariano Gálvez\\
Seminario de Tecnologías de Información\\
\textit{rramirezg18@miumg.edu.gt}
\end{center}

\vspace{10pt}


\subtitulo{Resumen}
El presente ensayo es acerca de las buenas practicas para realizar una encuesta efectiva y óptima, lo cual es importante para obtener información confiable que sirva de base para la toma de decisiones. Se investigó sobre la validez y la confiabilidad de los instrumentos de recolección de datos, además de que se abordó el diseño de las preguntas y de las escalas de respuesta como la escala Likert. Se describen los pasos que se deben seguir para planificar una encuesta desde la definición del objetivo hasta la distribución del instrumento y también se explican los sesgos más comunes que pueden afectar las respuestas de los encuestados. También se abordó la diferencia entre las preguntas abiertas y cerradas, así como el uso de las preguntas filtro y de control dentro de un cuestionario. Se concluye que una encuesta efectiva depende del diseño cuidadoso de las preguntas, de la validez y confiabilidad del instrumento, y de que este se dirija correctamente a la población y muestra de estudio.

\noindent\textbf{Palabras clave:} encuestas, validez, confiabilidad, escala Likert, recolección de datos

\vspace{6pt}

\subtitulo{Desarrollo del tema}

La encuesta es una búsqueda sistemática de información en la que el investigador pregunta sobre los datos que desea obtener y posteriormente reúne esos datos individuales para obtener datos agregados, además de que a diferencia de otras técnicas como la entrevista en profundidad, la particularidad de la encuesta es que se realizan las mismas preguntas a todos los entrevistados, en el mismo orden y en una situación social similar, lo cual se conoce como estandarización (Díaz de Rada, 2005). Díaz de Rada (2005) explica que esta estandarización es el fundamento básico de la encuesta ya que coloca a todos los entrevistados en una misma situación social, donde se les proporcionan los mismos estímulos por medio de las preguntas y posteriormente se recogen sus reacciones, de modo que las diferencias en las respuestas se deben a los propios entrevistados y no a la forma en que se aplicó el instrumento.

Toda encuesta se dirige hacia una población determinada, que es el conjunto total de personas o casos sobre los cuales interesa obtener información, sin embargo, por razones de tiempo y de costo, no es viable encuestar a toda la población por lo que se selecciona una muestra, que es un subconjunto representativo de esa población. Grasso (2006) señala que trabajar con una muestra permite ahorrar recursos financieros, de personal y de tiempo, lo cual permite dedicar los recursos disponibles a mejorar el trabajo de campo y reducir los errores de observación. Díaz de Rada (2005) agrega que la utilización de métodos de muestreo permite extender las conclusiones obtenidas en la muestra hacia ámbitos más amplios, lo cual solo es válido si la muestra seleccionada realmente representa a la población de estudio.

Sobre el proceso de investigación mediante encuesta, Díaz de Rada (2005) señala que aunque distintos autores dividen este proceso en diferente número de etapas, la mayoría coincide en cinco grandes momentos que son la formulación del problema, la descripción de la investigación, la obtención de datos, el análisis e interpretación de los datos y la presentación de resultados, lo cual permite contextualizar en qué parte de la investigación se ubica el trabajo de campo.

En cuanto al diseño de las preguntas del cuestionario, Díaz de Rada (2005) distingue entre preguntas abiertas y preguntas cerradas según la libertad de respuesta que se le da al entrevistado. Las preguntas abiertas no presentan alternativas de respuesta y el entrevistado responde con su propio vocabulario, lo cual exige un mayor esfuerzo del encuestador porque debe transcribir la respuesta tal y como es expresada, mientras que las preguntas cerradas ya presentan las alternativas de respuesta y facilitan tanto la aplicación como el posterior análisis estadístico de los datos.

Uno de los aspectos más importantes de un instrumento de recolección de datos es que sea válido y confiable. La validez se refiere a que el instrumento realmente mida lo que se pretende medir mientras que la confiabilidad tiene que ver con la exactitud y la consistencia de las mediciones que se obtienen (Corral, 2009). Corral (2009) explica que existen tres tipos de validez que son la validez de contenido, la validez de constructo y la validez predictiva o de criterio, además de que para estimar la confiabilidad de un instrumento se puede aplicar una prueba piloto a un grupo pequeño de entre 14 y 30 personas que tengan características similares a la muestra que se quiere estudiar.

Dentro de los tipos de preguntas cerradas también existen las preguntas filtro y filtradas, las cuales se utilizan para seleccionar entrevistados que cumplan con ciertas características y dirigirlos a un conjunto específico de preguntas, evitando así que los consultados den opiniones sobre asuntos que desconocen (Díaz de Rada, 2005). También existen las preguntas de control, que se utilizan para contrastar la consistencia de la información que da el entrevistado, formulando de manera distinta un aspecto que ya se preguntó antes en el cuestionario (Díaz de Rada, 2005).

Para medir actitudes y opiniones de forma subjetiva, Díaz de Rada (2005) menciona las escalas numéricas y de valoración, en las cuales el entrevistado se posiciona respecto a distintas opciones de respuesta según su propia percepción y no según una realidad objetiva. Una de las escalas más utilizadas de este tipo es la escala Likert, la cual consiste en presentar afirmaciones y pedir al encuestado que indique su nivel de acuerdo o desacuerdo con cada una (Matas, 2018). Matas (2018) señala que el formato de esta escala puede afectar la calidad de los datos obtenidos ya que existen sesgos como la tendencia central, que ocurre cuando los encuestados evitan las opciones extremas y prefieren marcar la alternativa intermedia, además de que también existe el sesgo de deseabilidad social en el cual la persona responde lo que considera que es mejor visto en lugar de su opinión real. Sobre el número de alternativas de respuesta, Matas (2018) recomienda utilizar escalas de cinco opciones junto con una alternativa de "no tengo opinión" para los casos en los que el encuestado no tenga suficiente información sobre el tema que se le pregunta.

El orden de las preguntas dentro del cuestionario también influye en la calidad de las respuestas, por lo que se recomienda aplicar la técnica del embudo, la cual consiste en iniciar con preguntas generales y poco sensibles para luego avanzar hacia preguntas más específicas y personales (Yepes, 2024).

Finalmente, antes de distribuir una encuesta de forma definitiva es recomendable hacer una prueba piloto con un grupo reducido de personas para identificar si existen preguntas ambiguas o mal redactadas y así poder corregirlas antes de aplicar el instrumento a la muestra completa (Netquest Blog, 2024).

\subtitulo{Observaciones y comentarios}
Considero que muchas veces se subestima el tiempo que se debe invertir en el diseño de una encuesta y se pasa directamente a redactar las preguntas sin definir antes el objetivo real de lo que se quiere medir. También pude notar que el tema de los sesgos en las escalas de respuesta no se toma en cuenta la mayoría de las veces al crear encuestas de forma rápida en herramientas en línea, lo cual afecta la calidad de los datos que se obtienen.

\subtitulo{Conclusiones}
\begin{enumerate}[label=\arabic*., left=0pt, itemsep=2pt]
\item Una encuesta efectiva depende principalmente de la validez y la confiabilidad del instrumento utilizado, además de que estas se deben comprobar antes de aplicar la encuesta a la población de estudio.
\item La redacción de las preguntas debe ser clara y neutral y debe evitar preguntas ambiguas, ya que esto afecta directamente la calidad de las respuestas obtenidas.
\item Considero que que una buena redacción en las preguntas es importante porque influye en las 
respuestas que da el entrevistado y esto afecta los datos que se obtienen.
\end{enumerate}

\subtitulo{Referencias}

\hangindent=1cm \noindent Díaz de Rada, V. (2005). \textit{Manual del trabajo de campo en la encuesta}. CIS - Centro de Investigaciones Sociológicas. \url{https://elibro.net/es/ereader/umg/52013}

\vspace{4pt}
\hangindent=1cm \noindent Grasso, L. (2006). \textit{Encuestas: elementos para su diseño y análisis}. Editorial Brujas. \url{https://elibro.net/es/ereader/umg/77141}

\vspace{4pt}
\hangindent=1cm \noindent Corral, Y. (2009). Validez y confiabilidad de los instrumentos de investigación para la recolección de datos. \textit{Revista Ciencias de la Educación}, \textit{19}(33), 228-247. \url{http://servicio.bc.uc.edu.ve/educacion/revista/n33/art12.pdf}

\vspace{4pt}
\hangindent=1cm \noindent Matas, A. (2018). Diseño del formato de escalas tipo Likert: un estado de la cuestión. \textit{Revista Electrónica de Investigación Educativa}, \textit{20}(1), 38-47. \url{https://doi.org/10.24320/redie.2018.20.1.1347}

\vspace{4pt}
\hangindent=1cm \noindent Netquest Blog. (2024, 22 de marzo). \textit{Cómo hacer una encuesta: Guía en 10 pasos}. Netquest. \url{https://www.netquest.com/blog/como-hacer-una-encuesta-guia-en-10-pasos}

\vspace{4pt}
\hangindent=1cm \noindent Yepes, V. (2024, 11 de diciembre). \textit{Escalas Likert: una herramienta fundamental en la ingeniería de encuestas}. El blog de Víctor Yepes - UPV. \url{https://victoryepes.blogs.upv.es/2024/12/11/escalas-likert-una-herramienta-fundamental-en-la-ingenieria-de-encuestas/}

\vspace{14pt}
\noindent\small \textbf{Encuesta:} https://forms.cloud.microsoft/r/12jQLcJgZb \\
\noindent\small \textbf{Código fuente:} https://github.com/rramirezg18/Foro-Academico-1.git

\end{document}